\subsection{Primitive local densities and generic class amplitudes}
\label{sec:flag-generic-local-amplitude}

We work at a finite place of the original unramified CM extension.
Write \(F_v\) for the local field, \(E_v\) for its unramified quadratic
algebra (including the split case), \(O\) for the ring of integers
of \(E_v\), \(\pi\) for the uniformizer of \(F_v\), and \(q\) for
its residue cardinality. The involution is denoted by a bar.
All Haar volumes below are relative to the compact subgroups
specified in the definitions; in particular we never replace a
reciprocal stabilizer volume by the number of its orbits.

\begin{lemma}[Uniform primitive density]\label{lem:flag-primitive-density}
Let \(V\) be a nondegenerate hermitian space of rank \(n\) over
\(E_v\) having a self-dual lattice \(L_0\). Suppose \(n\ge2m+3\),
where \(m\ge1\). For an integral hermitian \(m\times m\) matrix
\(B\), let \(P_j(B)\) count the ordered \(m\)-tuples in
\(L_0/\pi^j L_0\) with Gram matrix \(B\bmod\pi^j\) and with
primitive reduction (independent columns in each residue-field
factor). Then
\[
 P_j(B)=q^{(j-1)(2nm-m^2)}P_1(B),\qquad
 \tfrac34\le \beta(B):=q^{-(2nm-m^2)}P_1(B)\le\tfrac54 .
\]
The assertion includes residue characteristic two and singular
reductions of \(B\).
\end{lemma}
\begin{proof}
At a primitive tuple \(X\), the derivative of the Gram map is
\(Y\mapsto X^*HY+Y^*HX\). The map \(Y\mapsto X^*HY\)
is onto because \(X\) is primitive and \(H\) unimodular.
Every hermitian matrix is \(Z+Z^*\): use the off-diagonal entries
and the surjective residue trace for the diagonal. This works also
in characteristic two and in the split algebra. Thus every solution
at depth \(j\) has exactly \(q^{2nm-m^2}\) lifts at depth \(j+1\).

For the inert residue field, choose the columns successively.
After \(i\) independent columns, their prescribed pairings impose
\(i\) independent linear equations over \(\mathbb F_{q^2}\).
The direction space has dimension \(n-i\) and radical of dimension
at most \(i\). It has a nondegenerate complement of dimension
\(h\ge n-2i\). With the other coordinates fixed, completing the
square on this complement reduces the next norm condition to a
norm equation in dimension \(h\). The number of solutions differs
from \(q^{2h-1}\) by at most \((q-1)q^{h-1}\): indeed for a
nontrivial additive character the one-variable norm sum is \(-q\),
since every nonzero norm has \(q+1\) preimages, and Fourier
inversion gives the claimed bound in dimension \(h\). Subtracting
at most \(q^{2i}\) dependent vectors, the ratio of the number of
permitted choices to \(q^{2(n-i)-1}\) lies between
\[
 1-q^{1-(n-2i)}-q^{1-2(n-2i)},\qquad
 1+q^{1-(n-2i)}.
\]
The sum of the first errors over \(0\le i<m\) is at most
\(q^{-4}/(1-q^{-2})\le1/12\), and that of the second is
at most \(q^{-9}/(1-q^{-4})\le1/480\).
Multiplying the bounds gives the result (use
\(\prod(1-u_i)\ge1-\sum u_i\) and \(\prod(1+u_i)\le e^{\sum u_i}\)).

In the split case the two component matrices \(A,C\) have size
\(n\times m\), and the prescribed Gram condition is \(C^tA=B_1\).
First choose \(A\) of full rank, then choose the columns of \(C\)
in affine spaces of dimension \(n-m\), excluding the span of
previous columns. The normalized count is between
\[
 \prod_{i=0}^{m-1}(1-q^{i-n})(1-q^{i-(n-m)})\quad\hbox{and}\quad1 .
\]
The sum of the errors is at most
\((q^{-(n-m+1)}+q^{-(n-2m+1)})/(1-q^{-1})\le3/16\).
This proves the two inequalities.
\end{proof}

\begin{proposition}[Exact primitive-saturation mass]\label{prop:flag-primitive-mass}
Let \(J\subset V\) be a nondegenerate subspace of rank \(m\) and
let \(H_J=U(J^\perp)\) act as the identity on \(J\).
Fix its Haar measure, with the original reference compact subgroup
of volume one. For any full integral lattice \(M\subset J\), define
\[
 {\cal M}(M)=
 \sum_{\substack{\Lambda=\Lambda^\#,\ \Lambda\cap J=M\\
                   \bmod H_J(F_v)}}
 \frac1{\operatorname{vol}(H_J(F_v)\cap\operatorname{Aut}\Lambda)}.
\]
If \(x\) is a fixed basis of \(J\), and \(M\) has basis \(xA\), then
there is a single \(c_x>0\), independent of \(A\), such that
\[
 {\cal M}(M)=c_x|\operatorname{N}_{E_v/F_v}\det A|_v^{-(n-m)}
                      \beta(\operatorname{Gram}(xA)).
 \tag{GL1}
\]
In particular, in the range of Lemma~\ref{lem:flag-primitive-density},
every such integral \(M\) occurs. If a family of base lattices
has bases related by determinant-one changes, one can use the
same constant for the whole family. For an integral overlattice
\(M'\supset M\), its mass lies between
\(\tfrac34 c_x [M':M]^{-(n-m)}\) and
\(\tfrac54 c_x [M':M]^{-(n-m)}\) when \(x\) is a basis of \(M\).
\end{proposition}
\begin{proof}
Use additive Haar measures on tuples and hermitian matrices
normalized on \(L_0^m\) and on integral hermitian matrices.
The Gram map is a submersion at independent tuples: the same
derivative argument over \(F_v\) applies. Its quotient differential
measure on the nonsingular Gram fibre is invariant under \(U(V)\).
On each primitive residue cell, take the Gram coordinates and
\(2nm-m^2\) complementary coordinates. Their derivative is a
unit matrix up to integral changes of basis. Hensel lifting
(or the inverse function theorem with unit Jacobian on the
residue ball) identifies the cell with the product of these
residue balls. The measure on the fibre in the cell is therefore
\(q^{-(2nm-m^2)}\). Summing the primitive cells gives exactly
\(\beta\), not just a lower bound for a neighbourhood of the fibre.

By hermitian Witt extension, the fibre with nonsingular Gram is
the homogeneous space \(U(V)/H_J\) after choosing the given basis.
In the split case this is the corresponding statement about a
pair of dual vector spaces, obtained by extending bases of the
paired subspaces. The invariant fibre measure is a positive
multiple of quotient Haar measure. All self-dual lattices of \(V\)
are isometric: over the unramified field, a unimodular hermitian
form splits into unit norm lines by choosing a vector of unit norm
modulo \(\pi\), and unit norms are norms of units (residue norm
surjectivity and successive trace lifting). In the split case
a self-dual lattice is a lattice and its dual. Thus
\(K=U(L_0)\) gives all such lattices by \(U(V)/K\).
The double cosets for \(H_J\backslash U(V)/K\) with saturation \(M\)
correspond, by inversion, to the \(K\)-orbits on the primitive
Gram fibre in \(L_0^m\). A \(K\)-orbit has quotient volume
\(\operatorname{vol}(K)/\operatorname{vol}(H_J\cap gKg^{-1})\),
by quotient integration. Both groups are unimodular.
This gives the reciprocal stabilizer masses in the statement,
with the same choice of Haar measure on \(H_J\).

Finally right multiplication of tuples by \(A\) has ambient
Jacobian \(|\operatorname{N}\det A|_v^n\).
Congruence \(B\mapsto A^*BA\) on hermitian matrices has Jacobian
\(|\operatorname{N}\det A|_v^m\): check diagonal matrices
(on diagonal and paired off-diagonal coordinates) and elementary
matrices, which generate \(\mathrm{GL}_m(E_v)\). It follows that
the fibre measure is multiplied by
\(|\operatorname{N}\det A|_v^{n-m}\), whereas the homogeneous
quotient measure has not changed. This proves (GL1).
For an overlattice the displayed norm is its additive index,
as follows from elementary divisors in the field or split
factors. The count in the preceding lemma is positive, so the
primitive subset of the fibre is nonempty.
\end{proof}

\paragraph{The actual local frame classes.}
Now take the original two frames \(f,f'\) of length \(s=r+1\),
with their first \(r\) vectors isotropic and last-vector norms
\(-T,-T'\). Their cross matrix \(Z\) is the signed block-antidiagonal
point matrix, with consecutive block sizes \(s_1,\ldots,s_b\).
Inside a block its nonzero entries have a common scalar up to
sign. The original nonzero Gram invariant is
\[
 G=D_s\bar D_s-TT'D_{s-1}\bar D_{s-1},
 \qquad D_l=\det Z_{[1,l],[1,l]}.
\]
We assume the original local admission conditions, including
\(G\ne0\) and \(\theta=1\) or \(\pi\), and the existence of an
admitted term. Put \(m=2s\), \(n=k+2r+1\), and
\[
 a=v(G)-m v(\theta),\quad
 d=\sum_{D_l\ne0,\ l>0}v(D_l\bar D_l),\quad
 P=q^{v(\theta)+a+d}.
 \tag{GL2}
\]
These are the source's \(a,vD,Hp\). The corrected joint lattice
of a pair of adapted frames \(u,u'\) is
\[
 M(u,u')=\langle u_i,\theta^{-1}(u_i'-u_i)\rangle_O .
\]
The determinant of its Gram matrix is
\((-1)^s\theta^{-2s}G\), by the joint Gram determinant identity
and the indicated change of basis. Its span \(J\) is nondegenerate.
For an admitted pair the corrected generators are integral,
and every endpoint minor is integral. Thus \(a,d\ge0\);
\([M^\#:M]=q^{2a}\). This also gives the defect bound:
for every integral overlattice \(M_h\) of \(M\) in \(J\),
\[
 M\subset M_h\subset M_h^\#\subset M^\#,\qquad
 [M_h:M]\le q^a.
 \tag{GL3}
\]
The number of such overlattices is at most \(q^{2am}\).
Indeed every submodule of \(M^\#/M\) has at most \(m\) generators
over \(O\), by the elementary divisor theorem (choose generators
componentwise in the split case). Count all ordered lists of
\(m\) elements. In particular \(P=1\) forces \(M=M_h\).

Let \(U\) denote the upper unitriangular group of size \(s\).
Write \(S=N_f\cap N_{f'}\) for the rational common unipotent
stabilizer over \(F_v\), and \(H_J=U(J^\perp)\).
The two actions commute. A common unipotent is determined by
its two adapted frame matrices \(A,B\in U(E_v)\), with
\(A^*ZB=Z\). It fixes \(J^\perp\) pointwise since the difference
on that subspace lies in the intersection of the two isotropic
spans, which is zero. Conversely, a pair satisfying this equation
defines an isometry of \(J\) and extends by the identity on
\(J^\perp\). It belongs to both flag unipotents: it acts trivially
on the successive isotropic flag quotients and on the middle
quotients. For the last assertion, the latter quotient inside
\(J\) is spanned by the last vector of the first frame and
the line in the second frame orthogonal to the first \(r\)
vectors; the last vector is fixed modulo the isotropic span
and the line is fixed by the unipotent action. The same
argument applies in the other order and separately in the
two split factors.

As \(Z\) reverses each consecutive block, conjugating an upper
triangular matrix by \(Z\) sends its upper entries to lower ones
exactly within the same block. Hence \(B\) must be block diagonal,
and \(A\) is determined by \(B\). The group is a product of
upper unitriangular groups of sizes \(s_j\); the within-block
signs do not affect integrality. Its reference compact
\(K_0\) is precisely integral \(B\) in these blocks.
Normalize \(\operatorname{vol}(K_0)=1\).

The frame-class quotient is the double quotient
\[
 S\backslash (U(E_v)\times U(E_v))/Q_\theta,\qquad
 Q_\theta=\{(A,B)\in U(O)^2:A\equiv B\bmod\theta O\}.
\]
An adapted change of each frame is realized by its flag
unipotent, by extending a triangular change on the isotropic
vectors and a shift of the last vector into the isotropic
span. For clarity the extension can be made with a hyperbolic
splitting: fix the isotropic vectors \(e_i\), add \(h_i\) in the
middle space to their duals \(d_i\), and send a middle vector
\(w\) to \(w-\sum_i h(w,h_i)e_i\) (use the convention linear
in the first variable). Add an isotropic-coordinate matrix to
the images of the \(d_i\) whose sum with its adjoint is minus
their middle Gram. The trace-one identity supplies the diagonal
even at dyadic places. The arbitrary triangular change on
the \(e_i\) is extended by its inverse adjoint on the duals.
This construction also gives integral extensions when the
hyperbolic splitting and the shift coefficients are integral
and the last vector is primitive in its self-dual middle lattice.

An integral frame class is one with \(Y=A^*ZB\) integral.
Its weight is
\[
 w(C)=1/\operatorname{vol}(S\cap pQ_\theta p^{-1}),\qquad
 p=(A,B).
\]
This is exactly the original \(wC\) since that intersection
preserves the linked integral frame class. We next prove
\[
 w(C)\ge1,\qquad
 \sum_{C\ {\rm integral}}w(C)\le
       q^{s(s-1)(d+v(\theta))}\le P^{s(s-1)}.
 \tag{GL4}
\]

Normalize the representative in each left \(S\)-coset by
requiring the diagonal blocks of \(B\) to be identity.
This uniquely fixes the \(S\)-action. Its remaining coordinates
are those of the lower unitriangular \(L=A^*\) and off-block
entries of \(B\); there are at most \(s(s-1)\) coordinates
over \(E_v\). Haar disintegration in these coordinates
is ordinary product additive measure: the block-diagonal
part of \(B\) is the \(S\)-coordinate, and the other changes
are multiplication by unitriangular matrices, with triangular
Jacobian of absolute value one. Each \(O\) has additive
volume one. At the normalized representative the block-diagonal
part of the second component of \(p^{-1}Sp\) is unchanged.
Thus its membership in \(Q_\theta\) forces the
\(S\)-coordinate to be integral; the stabilizer is contained
in \(K_0\), proving \(w(C)\ge1\).

Here are explicit denominator bounds in the normalized
coordinates. The leading block determinants of \(Y=LZB\)
are the corresponding \(D_l\). For a column block \(j>i\),
the last block row of the inverse leading square of \(Y\)
through block \(i\), applied to the leading rows of column
block \(j\), is \(B_{ij}\). Its entries have denominators
dividing \(D_{\mathrm{end}(i)}\). The Schur complement of the
leading square through block \(i-1\), evaluated in row
block \(j\ge i\) and column block \(i\), equals
\(L_{ji} Z_i\). Therefore the entries of \(L_{ji}\) have
denominators dividing \(D_{\mathrm{end}(i-1)}z_i\), where
\(z_i\) is the scalar of \(Z_i\) and \(D_0=1\).
At each factor of \(E_v\) the valuation of this denominator is
\[
 (1-1/s_i)v(D_{\mathrm{end}(i-1)})
             +(1/s_i)v(D_{\mathrm{end}(i)}).
\]
It is nonnegative, and its norm valuation is at most \(d\).
The same bound holds for the denominators of \(B_{ij}\).
Consequently each coordinate lies in a fixed fractional ideal
of additive volume at most \(q^d\). The entire integral region
in the left quotient has volume at most \(q^{s(s-1)d}\).
The compact subgroup \(Q_\theta\) has volume
\(q^{-s(s-1)v(\theta)}\) in the product measure.
Quotient integration assigns to each right \(Q_\theta\)-orbit
the volume \(\operatorname{vol}(Q_\theta)w(C)\).
This proves (GL4), including finiteness. If \(P=1\) the
normalized coordinates are all integral and \(\theta=1\);
there is exactly one occupied class and its weight is one.

\paragraph{The datum fibres and their amplitudes.}
A local datum consists of the original two inner self-dual
quotient lattices, with primitive last vectors, and their
level graph. Use the datum and term definitions in
\texttt{LocNG} or equivalently their \texttt{Frames} coordinates.
For a fixed datum \(\delta\) and frame class \(C\), let
\(N_\delta(C)\) count actual lattices in that class, without
quotienting by the common unipotent. Write \(N_\delta^0(C)\)
for those satisfying the generic condition
\(\Lambda\cap J=M(u,u')\), with \(u,u'\) any representative
of \(C\). The lattice \(M(u,u')\) is independent of the
representative since the linked changes preserve its generators.

We detail the witness dictionary. In a self-dual lattice with
primitive isotropic vectors and primitive last vector modulo
their span, choose duals to the whole primitive frame using
unimodularity. In particular the isotropic duals pair to zero
with the last vector. Correct their mutual Gram by isotropic
coordinates using the surjective integral trace. This splits
off the hyperbolic planes and gives a self-dual middle lattice.
The rational isometry from the reference hyperbolic splitting
to this one, inducing the identity on the quotient by the
isotropic span, lies in the indicated flag unipotent. This
constructs the two witnesses in the local term definition
from such a lattice and its two adapted primitive frames.
Conversely those witnesses give the stated primitive frames.
For two witnesses of the same lattice, their adapted frames
differ by integral triangular matrices: each is a basis
of the same primitive frame sublattice with fixed graded
coordinates. Since both pairs satisfy
\(u_i'\equiv u_i\bmod\theta\Lambda\), these matrices are
congruent modulo \(\theta\) by primitivity.
Conversely every linked integral change is realized in the
same lattice by the integral unipotent extension above.
The quotient lattices and the level graph are intrinsic to
\(\Lambda\), since the unipotents act trivially on the
two quotients. This proves the frame-class dictionary.

If \(\Lambda\cap J=M(u,u')\), both frames are primitive:
each is part of a basis of \(M\) (use the reverse
corrected generators for the second), and \(M\) is saturated.
Thus every self-dual lattice of this saturation gives a
generic term, with its induced datum. For arbitrary integral
overlattices of \(M\) we need only the inclusion of admitted
terms into self-dual lattices of that saturation.

Fix a datum \(\delta\) and an integral class \(C\).
The common complement lattice \(\Lambda\cap J^\perp\)
of any generic term is determined by \(\delta\).
Indeed for \(w\in J^\perp\), membership in its two inner
lattices together with the level graph means that there are
lifts \(w+x\), \(w+x'\) in \(\Lambda\), where \(x,x'\) lie in the
respective isotropic spans, and \(x'-x\in\theta\Lambda\).
Away from the level this last condition is automatic.
By generic saturation,
\(\theta^{-1}(x'-x)\in M(u,u')\). Comparing coefficients
in the independent combined frame, both \(x\) and \(x'\)
are integral combinations of the respective frame vectors.
It follows that \(w\in\Lambda\); the converse is immediate.
Call this fixed complement lattice \(L\).

For prescribed saturated lattices \(M,L\) in the orthogonal
summands, a self-dual lattice of \(V\) with these intersections
is the graph of an anti-isometry of the finite hermitian
discriminant modules \(M^\#/M\) and \(L^\#/L\).
Indeed its image in their direct sum is self-orthogonal
and equal to its orthogonal complement, and neither coordinate
projection has a kernel by saturation; the order equality
then forces both projections to be bijective.
This gives a bijection with anti-isometries, if any exist.
There are at most \(|M^\#/M|^m=q^{2am}\) such maps, since
the first module is generated by at most \(m\) elements.
Thus
\[
 N_\delta^0(C)\le q^{2am}. \tag{GL5}
\]

We now account for the two stabilizer weights. The unipotent
\(S\) preserves the datum, and the stabilizer of \(C\) is a
compact open subgroup \(S_C\). Its orbits on lattices in
that class have cardinality
\(\operatorname{vol}(S_C)/\operatorname{vol}(S_\Lambda)\).
Consequently, with \(\mu(\Lambda)=1/\operatorname{vol}(S_\Lambda)\),
the sum of \(\mu\) over the unipotent orbits in the datum and
class is exactly \(w(C)N_\delta(C)\); likewise generically.
The complement group fixes each pair of frames and acts
equivariantly on the datum. Give it the Haar measure in which
the reference datum stabilizer has volume one; thus
\(\mathrm{wt}(\delta)=1/\operatorname{vol}(H_\delta)\).
These stabilizers are compact open: they preserve full
quotient lattices, whose intersections with the complement
give full lattices there, and the additional level graph
is finite residue data. For each complement orbit of
data, the complement-orbits of its lattices are the
\(H_\delta\)-orbits of the fixed fibre. The same orbit-size
calculation gives the identities and inequality
\[
 \sum_{\delta\bmod H_J}\mathrm{wt}(\delta)N_\delta^0(C)
     ={\cal M}(M(u,u')),\qquad
 \sum_{\delta\bmod H_J}\mathrm{wt}(\delta)N_\delta(C)
     \le\sum_{M_h\supset M(u,u')}{\cal M}(M_h).
 \tag{GL6}
\]
These calculations are valid for nonnegative sums before
finiteness has been established: lattices form a countable
set and all compact stabilizer orbits on them are finite.
No assumption of a unique integral lattice in a rational
double coset has been made.

Every local phase in the amplitude has modulus one.
Changing witnesses by the linked integral frame changes does
not change the phase: the difference of the two simple-root
sums is divisible by \(\theta\), and the underlying additive
character is integral-trivial. Common unipotent translation
cancels by relevance. Hence the phase is the original
well-defined phase on the indicated lattice orbits.

\begin{theorem}[Local generic amplitude and mass]\label{thm:flag-generic-local-amplitude}
In the original rank range \(k\ge4r+7\), at every occupied
local nonzero-Gram point set \(e=s(s-1)+4s\).
Let \(m_\delta,m_\delta^0\) be respectively the total and
generic unipotent mass in the datum \(\delta\), and \(A_\delta^0\)
the generic amplitude, with its original phase.
Let \(nn=\sum_\delta\mathrm{wt}(\delta)m_\delta\).
For any occupied integral pairing class \(C_*\), let
\(ngen(C_*)\) be the generic mass in that class, summed over
data with the same weights. Then
\[
 \sum_\delta\mathrm{wt}(\delta)|A_\delta^0|^2
       \le P^e nn,\qquad
 0<ngen(C_*)<\infty,\qquad nn\le2P^e ngen(C_*).
 \tag{GL7}
\]
At \(P=1\) the two comparisons are equalities (with factor
one, not two). These give the generic local bound with
\(C_1=39\), and the local pairing-class lower bound with
\(C_0=9\), in the original concrete finite-place setting.
\end{theorem}
\begin{proof}
By (GL4)--(GL5), \(m_\delta^0\le P^e\).
Thus \(|A_\delta^0|^2\le(m_\delta^0)^2\le P^e m_\delta\).
For the second bound use (GL6), (GL3) and (GL1).
All corrected base frames are related by determinant-one
changes of the corrected basis, even before integrality,
so the constant \(c_x\) is the same for every class.
Here \(n-m=k-1\) and \(n\ge2m+3\).
Each class contributes at most \(\tfrac54 c_x q^{2am}w(C)\)
to \(nn\). The chosen occupied class contributes
\[
 ngen(C_*)=w(C_*){\cal M}(M(C_*))
                  \ge\tfrac34 c_x
\]
generically; it is finite by (GL1). Thus \(nn\le
(5/3)P^e ngen(C_*)\le2P^e ngen(C_*)\).

To see that this is the source pairing class, two pairs
with cross matrices related by linked integral changes
can first be so adjusted to have equal cross matrices.
The resulting isometry on \(J\), fixing the complement
pointwise, is in the two unipotents by the frame argument
above. Thus equality of the pairing class is precisely
equality of the double frame class modulo \(S\).
There is no additional factor from its individual
integral frame representatives.

At \(P=1\) there is one class of weight one, every saturation
is generic, and (GL5) gives at most one lattice per datum
in a fixed class. Every nonempty datum has mass one and
a single phase of modulus one. Therefore the amplitude square,
the total mass and the generic mass coincide. Finally
\(e=s(s-1)+4s=s^2+3s\le39s^2\) and
\(e\le9s^2+4s\), also at \(r=0\).
\end{proof}

The proof of (GL7) uses no cancellation among nongeneric
terms; it is the generic part of the local input, not the
nongeneric saving \(q^{-(k-2r-5)}\).
The global class unfolding and the comparison with the
pinned real orbital integral are separate obligations.
